% ------------------------------------------------------------
% Quenya-/Elbische FDI-Schriften
% ------------------------------------------------------------
% Wird von \FDIInputLanguageFonts in der Praeambel geladen.
% Wichtig: Diese Fonts werden NICHT global auf das Dokument angewendet,
% sondern nur ueber die FDI-Font-Makros verwendet.
% Renderer=Basic ist bei alten Tengwar-/Symbolfonts oft wichtig.
%
% Der Pfad muss hier ausgeschrieben werden, weil \newfontfamily den
% Pfad schon beim Laden der Praeambel aufloest und \FDILanguage zu
% diesem Zeitpunkt noch die primaere Comic-Sprache ist.

% Normale FDI-Sprechblasen / Standard-Overlay
\newfontfamily{\FDIQyaArialFont}{tngani_qya.ttf}[
  Path = language_files/qya/fonts_qya/,
  Renderer = Basic,
  Scale = 0.78
]

% Fette FDI-Sprechblasen / Hervorhebungen
\newfontfamily{\FDIQyaArialBOLDFont}{tngani_qya.ttf}[
  Path = language_files/qya/fonts_qya/,
  Renderer = Basic,
  Scale = 0.78,
  FakeBold = 1.8
]

% FDI-Overlay-Schrift, etwas groesser fuer Comic-Texte
\newfontfamily{\FDIQyaOverlayFont}{tngani_qya.ttf}[
  Path = language_files/qya/fonts_qya/,
  Renderer = Basic,
  Scale = 0.82
]

% FDI-Karolingische-Minuskel-Ersatzschrift
\newfontfamily{\FDIQyaKarolingischeMinuskel}{tngani_qya.ttf}[
  Path = language_files/qya/fonts_qya/,
  Renderer = Basic,
  Scale = 0.86,
  FakeBold = 1.3
]

% Zuweisung der FDI-Font-Makros fuer diese Sprache.
% Wird von \FDISetupLanguageFonts ausgefuehrt, sobald qya gesetzt ist.
\FDISetLanguageFonts{qya}{%
  \let\KarolingischeMinuskel\FDIQyaKarolingischeMinuskel
  \let\FDIOverlayFont\FDIQyaOverlayFont
  \let\ArialBOLDFont\FDIQyaArialBOLDFont
  \let\ArialFont\FDIQyaArialFont
}
